% paper/sections/10_3_level_selection.tex
%
% §10.3  Level 1/2/3 選択ロジック（SP-I §3）
%
% ═══════════════════════════════════════════════════════════════════════════════
\subsection{Level 1/2/3 選択ロジックとコスト-精度トレードオフ}
\label{sec:level_selection}
% ═══════════════════════════════════════════════════════════════════════════════

SP-I \cite{SPI} §3 で時間積分設計を Level 1--3 の連続スペクトルに整理した
（§\ref{sec:time_level1}--\ref{sec:time_level3}）．
本節では 1 タイムステップ内で Level を切替えるトリガー，
Level 間コスト比率，および適用領域の棲み分けを示す．

% -------------------------------------------------------------------------------
\subsubsection{Level 選択トリガー：剛性レジームの自動判定}
\label{sec:level_trigger}
% -------------------------------------------------------------------------------

物理 CFL 制約の 4 つを通常ステップで測定し，最も厳しい制約に対応する
Level を選ぶ：
\begin{equation}
  \Delta t_\text{step}
    = \min\!\Bigl(
        \Delta t_\text{adv},\,
        \Delta t_\text{visc},\,
        \Delta t_\sigma^\text{BKZ},\,
        \Delta t_\text{wave}
      \Bigr),
  \label{eq:level_cfl_min}
\end{equation}
ここで
\begin{itemize}[noitemsep]
  \item $\Delta t_\text{adv} = C_\text{adv} h/\|\mathbf{u}\|_\infty$（移流 CFL, $C_\text{adv}\le 0.5$）
  \item $\Delta t_\text{visc} = C_\text{visc}\rho h^2/(2\mu N_\text{dim})$（粘性 CFL）
  \item $\Delta t_\sigma^\text{BKZ} = \sqrt{(\rho_l+\rho_g)h^3/(2\pi\sigma)}$
        （Brackbill--Kothe--Zemach 毛管波制約，§\ref{sec:time_capillary}）
  \item $\Delta t_\text{wave} = C_\text{wave}h/c_s$（物理波速度，非圧縮では無関係）
\end{itemize}

\begin{table}[htbp]
\centering
\small
\begin{tabular}{@{}lllll@{}}
\hline
最少制約 & 支配剛性 & 推奨 Level & 対応スキーム & 典型領域 \\
\hline
$\Delta t_\text{adv}$ & 対流 &
  Level 1 / 2 & SSPRK3 / AB2 &
  一般流れ，Re 中庸 \\
$\Delta t_\text{visc}$ & 粘性 &
  Level 2 & AB2 + CN &
  高粘性比，薄膜 \\
$\Delta t_\sigma^\text{BKZ}$ & 毛管波 &
  Level 2 / 3 & 半陰的 ST / Li 2022 &
  We $\ll 1$，微小気泡 \\
同時多数 & 結合剛性 &
  Level 3 & Radau IIA + JFNK &
  高密度比 + 低 We \\
\hline
\end{tabular}
\caption{CFL 最少制約から Level を自動選択する規則．
動的切替は実行時 $\max$ の統計を 10 ステップごとに監視し，
Level が連続 3 回推奨と一致したときのみ切替える（チャタリング防止）．}
\label{tab:level_trigger}
\end{table}

% -------------------------------------------------------------------------------
\subsubsection{Level 間コスト-精度トレードオフ}
\label{sec:level_cost_accuracy}
% -------------------------------------------------------------------------------

1 タイムステップ当たりのコストを\textbf{PPE 求解 1 回分を 1 単位}として比較する：

\begin{center}
\begin{tabular}{@{}lllll@{}}
\hline
Level & 1 ステップコスト & 許容 $\Delta t$ & 実効 throughput & 切断誤差 \\
\hline
Level 1 & $\sim$ 3--5 (RK 段) &
  $\min(\text{CFL})$ & 1.0× (ベース) &
  $\Ord{\Delta t^2}$ 時間，空間は演算子依存 \\
Level 2 & $\sim$ 6--10 (AB2+CN+ST) &
  $10$--$50\times$ Level 1 & $2$--$8\times$ Level 1 &
  $\Ord{\Delta t^2}$ 全体 \\
Level 3 & $\sim$ 30--100 (Radau IIA + JFNK) &
  $100$--$1000\times$ Level 1 & $3$--$10\times$ Level 2 &
  $\Ord{\Delta t^5}$ 全体 \\
\hline
\end{tabular}
\end{center}

\noindent\textbf{解釈：}
Level 2 は $C_\text{visc}$ と $C_\sigma$ の両制約を同時に緩和するため，
ch13 水--空気級の毛管駆動シミュレーションでは
Level 1 の 8 倍程度の throughput を実現する．
Level 3 は外殻非線形反復のコストが高いが，
$\Delta t$ 拡大比が剛性の平方で効くため，
$\rho_l/\rho_g > 500$ 領域で Level 2 を上回る実効効率を持つ．

% -------------------------------------------------------------------------------
\subsubsection{バリデーション Level: Level 1 の役割}
\label{sec:level1_validation}
% -------------------------------------------------------------------------------

Level 1（SSPRK3 全陽的）は\textbf{実運用ではなく}検証用途として
以下の場面に限定する：
\begin{itemize}[noitemsep]
  \item グリッド収束試験（§\ref{sec:verify_ccd_convergence}, §\ref{sec:verify_ppe_neumann}）：
        時間離散化誤差を空間誤差の 1/10 以下に抑え
        空間精度 $\Ord{h^k}$ を純粋に測定．
  \item コンポーネント単体テスト：
        個別演算子（FCCD 勾配，UCCD6 対流，CCD 曲率）の整合性を SSPRK3 で確認．
  \item H-01 診断（§\ref{sec:h01_diagnosis_fccd_remedy}）：
        BF 残差の切断誤差次数を Level 2 の半陰的結合項抜きで測定する．
\end{itemize}

Level 1 で得た精度指標を\textbf{Level 2 の設計目標}として引き継ぐ
（SP-I §5, \cite{SPI}）．
Level 2 の Crank--Nicolson + 半陰的 ST は Level 1 の $\Ord{h^k}$ を
時間刻み制約から切り離した上で保持することを設計要件とする．

% -------------------------------------------------------------------------------
\subsubsection{プロダクション Level: Level 2 の既定}
\label{sec:level2_production}
% -------------------------------------------------------------------------------

Level 2（AB2 + CN + 半陰的 ST + BF 投影）は本稿の\textbf{既定プロダクション構成}．
§\ref{sec:time_level2} で定式化した通り，時間刻みは
\begin{equation}
  \Delta t_\text{L2}
    = \min\!\bigl(
        C_\text{adv} h/\|\mathbf{u}\|_\infty,\,
        C_\sigma^\text{impl}\,\Delta t_\sigma^\text{BKZ}
      \bigr),
  \qquad C_\sigma^\text{impl}\simeq 5\text{--}10,
  \label{eq:level2_dt}
\end{equation}
半陰的 ST が $\Delta t_\sigma$ の実用係数を $\sim 10\times$ 緩和するため，
移流 CFL のみが実効制約となる．
境界条件設定（Option III/IV, §\ref{sec:fccd_bc_options}），
BF 七原則（P-1..P-7, §\ref{sec:bf_seven_principles}），
面ジェット統一（§\ref{sec:fccd_bf_sub_system}）は
Level 2 で全て要求される．

% -------------------------------------------------------------------------------
\subsubsection{極端剛性 Level: Level 3 の適用窓}
\label{sec:level3_stiff_regime}
% -------------------------------------------------------------------------------

Level 3（Radau IIA + JFNK + 完全陰的 ST）は以下の\textbf{狭い適用窓}で
Level 2 を上回る：
\begin{itemize}[noitemsep]
  \item $\rho_l/\rho_g \ge 500$ かつ $\mu_l/\mu_g \ge 50$ かつ $\text{We} \le 0.01$
        の同時成立．
  \item 薄膜流，マイクロ流体，インクジェット breakup，
        合体・分裂開始時の界面極限挙動．
  \item $\Delta t_\sigma^\text{BKZ}/\Delta t_\text{adv} < 0.01$
        （毛管波が対流を凌駕）．
\end{itemize}

Level 3 の外殻 JFNK 反復は 1 ステップあたり 5--20 Krylov 反復 + 3--6 Newton
ステップを要するが，$\Delta t_\text{L3}/\Delta t_\text{L2}$ が
$50\text{--}200\times$ まで拡大可能なため実効効率が上回る．
本稿では Level 3 の具体構成（Radau IIA の stage 構成，JFNK 前処理戦略）は
§\ref{sec:time_level3} にとどめ，詳細実装は SP-I §3.3 \cite{SPI} を参照とする．

% -------------------------------------------------------------------------------
\subsubsection{Level 選択マトリクス（物理 → 推奨）}
\label{sec:level_physics_matrix}
% -------------------------------------------------------------------------------

\begin{center}
\small
\begin{tabular}{@{}lllll@{}}
\hline
物理例 & $\rho_l/\rho_g$ & $\text{We}$ & 推奨 & 理由 \\
\hline
2 次元渦流（TGV）& $1$ & $\infty$ & L1/L2 & 検証用，移流 CFL のみ \\
上昇気泡（中程度）& $100$ & $1$ & L2 & BF + 半陰的 ST が実用 \\
毛管波（ch13）& $830$ & $0.02$ & L2 & $\Delta t_\sigma^\text{BKZ}$ 緩和が鍵 \\
水--空気インクジェット & $830$ & $0.001$ & L2/L3 & 破断初期は L3 必要 \\
マイクロ流体 T 字路 & $1000$ & $0.0001$ & L3 & 全剛性同時，L2 では step が停止 \\
Zalesak 型（硬幾何）& $\{1\}$ & N/A & L1 & CLS 単独検証 \\
\hline
\end{tabular}
\end{center}

実装上は YAML 設定で
\verb|time_integration.level: 1|\,\verb|2|\,\verb|3|
を明示指定し，ランタイムトリガー（表~\ref{tab:level_trigger}）は
警告のみ発行する．
Level の自動昇格は収束性と再現性を損なうため既定では無効化する
（SP-I §6, \cite{SPI}）．

% ═══════════════════════════════════════════════════════════════════════════════
